\documentclass[11pt]{article}
\usepackage{mathblatt}
% vorspann.tex – Kompetenzblatt, zusammenbau v0.7 (2026-09-28).
% Ergänzt mathblatt.sty für Kompetenzblätter, ohne sie zu ändern
% (layout-befunde.md 1–34 vom 28.09.). Wird nach \usepackage{mathblatt}
% eingelesen.
\makeatletter
\RequirePackage{newunicodechar}
% Zeichen, die Latin Modern nicht hat (Befund 20): Text aus einer
% Ersatzschrift, Mathe als Befehl; ohne Ersatzschrift auch Text als Mathe.
\newif\ifkbersatz
\IfFontExistsTF{FreeSerif}{\newfontfamily\kbersatzschrift{FreeSerif}\kbersatztrue}{%
 \IfFontExistsTF{DejaVu Serif}{\newfontfamily\kbersatzschrift{DejaVu Serif}\kbersatztrue}{%
  \IfFontExistsTF{Cambria Math}{\newfontfamily\kbersatzschrift{Cambria Math}\kbersatztrue}{%
   \IfFontExistsTF{Segoe UI Symbol}{\newfontfamily\kbersatzschrift{Segoe UI Symbol}\kbersatztrue}{}}}}
\newcommand{\kbzeichen}[2]{\ifmmode #2\else\ifkbersatz{\kbersatzschrift #1}\else\ensuremath{#2}\fi\fi}
\newcommand{\kbzeichenhoch}[1]{\ifkbersatz\ifmmode\text{\kbersatzschrift #1}\else{\kbersatzschrift #1}\fi\else ?\fi}
% Kopf- und Fußzeile (Befund 1, 2): Kopf Thema · Blattart, Fuß Kennung und Seite
\newcommand{\kbkopfzeile}[2]{\pagestyle{fancy}\fancyhf{}\renewcommand{\headrulewidth}{0pt}%
  \fancyhead[L]{\footnotesize\color{mbgrau}#1}%
  \fancyfoot[L]{\footnotesize\color{mbgrau}#2}%
  \fancyfoot[R]{\footnotesize\color{mbgrau}Seite \thepage}}
% Flattersatz überall (Befund 25)
\raggedright
% Titel des Blatts: Ich-kann-Satz, darunter Zweigzeile
\newcommand{\kbtitel}[2]{\par\noindent{\Large\bfseries\raggedright #1\par}%
  \ifblank{#2}{}{\par\smallskip\noindent{\small #2}\par}\medskip}
% Abschnitt: „Das kennst du schon“, „Zum Merken“
\newcommand{\kbabschnitt}[1]{\par\addvspace{12pt}\Needspace*{5\baselineskip}%
  \noindent{\bfseries #1}\par\nobreak\vspace{1pt}\noindent{\color{mbgrau}\rule{\linewidth}{0.4pt}}\par\nobreak}
% Hauptnummer: Titel hängt am ersten Block; jeder Block (Teilaufgabe) bricht
% nicht, passt er nicht mehr, rückt er ganz auf die nächste Seite (Befund 21).
\newif\ifkbtitel
\newsavebox{\kbbox}
\newlength{\kbhoehe}
\newlength{\kbnummerabstand}\setlength{\kbnummerabstand}{10pt}
\newlength{\kbteilabstand}\setlength{\kbteilabstand}{6pt}
\newenvironment{kbaufgabe}[1]{\stepcounter{aufgabe}\setcounter{teil}{0}%
  \gdef\kbtiteltext{\noindent\hangindent1.6em\hangafter1\makebox[1.6em][l]{\textbf{\theaufgabe.}}#1\par}%
  \global\kbtiteltrue\par\addvspace{\kbnummerabstand}}{\par}
\newenvironment{kbblock}{\par\addvspace{\kbteilabstand}\begin{lrbox}{\kbbox}%
  \begin{minipage}[t]{\linewidth}\raggedright\parskip\z@
  \ifkbtitel\kbtiteltext\vspace{2pt}\global\kbtitelfalse\fi
  \list{}{\leftmargin1.6em\labelwidth1.4em\labelsep0.2em\topsep\z@\partopsep\z@
    \parsep\z@\itemsep\z@\rightmargin\z@}\raggedright}%
  {\endlist\end{minipage}\end{lrbox}%
  \setlength{\kbhoehe}{\dimexpr\ht\kbbox+\dp\kbbox\relax}%
  \Needspace*{\kbhoehe}\noindent\usebox{\kbbox}\par}
% Paarsatz (Platz nutzen, Befund 27): zwei Hauptnummern oder zwei
% Teilaufgaben mit Zeichenaufgabe nebeneinander, als ein Block
\newcommand{\kbnummer}[1]{\stepcounter{aufgabe}\setcounter{teil}{0}%
  \noindent\hangindent1.6em\hangafter1\makebox[1.6em][l]{\textbf{\theaufgabe.}}#1\par\vspace{2pt}}
\newcommand{\kbhalb}[1]{\list{}{\leftmargin1.6em\labelwidth1.4em\labelsep0.2em\topsep\z@
  \partopsep\z@\parsep\z@\itemsep\z@}\raggedright #1\endlist}
\newcommand{\kbpaar}[2]{\par\addvspace{\kbnummerabstand}\begin{lrbox}{\kbbox}%
  \begin{minipage}[t]{0.485\linewidth}\raggedright\parskip\z@ #1\end{minipage}\hfill
  \begin{minipage}[t]{0.485\linewidth}\raggedright\parskip\z@ #2\end{minipage}\end{lrbox}%
  \setlength{\kbhoehe}{\dimexpr\ht\kbbox+\dp\kbbox\relax}\Needspace*{\kbhoehe}\noindent\usebox{\kbbox}\par}
\newcommand{\kbteilpaar}[2]{\item[]\noindent
  \begin{minipage}[t]{0.485\linewidth}\kbhalb{#1}\end{minipage}\hfill
  \begin{minipage}[t]{0.485\linewidth}\kbhalb{#2}\end{minipage}\par}
% Anweisung der Hauptnummer, eigene Zeile über der Teilaufgabe (Befund 5)
\newcommand{\kbanweisung}[1]{\item[]#1\par\vspace{2pt}}
% Kopfzeile einer Teilaufgabe: Buchstabe, Auftakt (halbfett oder Kontext),
% Prüfkennung klein rechts (Befund 3, 12, Beschluss c)
\newlength{\kbkennbreite}\setlength{\kbkennbreite}{1.9cm}
\newcommand{\kbkennung}[1]{{\footnotesize #1}}
\newcommand{\kbteil}[3]{\item[#1]\ifblank{#3}%
  {\parbox[t]{\linewidth}{\raggedright #2}}%
  {\parbox[t]{\dimexpr\linewidth-\kbkennbreite\relax}{\raggedright #2}\hfill
   \parbox[t]{\kbkennbreite}{\raggedleft\kbkennung{#3}}}\par}
% Frage in eigener Zeile (Befund 4), ein Satz je Zeile (Befund 10)
\newcommand{\kbfrage}[1]{\par\vspace{2pt}#1\par}
% Antwortfeld in eigener Zeile darunter, linksbündig (Befund 4, 8)
\newcommand{\kbantwort}[1]{\par\vspace{4pt}\noindent#1\par}
% Zwei Spalten: links Grafik/Tabelle/Aufgabe, rechts Antwort oder Rechenraum
% (Befund 27); oben bündig. Beginnt einen eigenen Listenpunkt ohne Marke.
\newcommand{\kbzweispaltig}[3][0.48]{\item[]\noindent
  \begin{minipage}[t]{#1\linewidth}\raggedright #2\end{minipage}\hfill
  \begin{minipage}[t]{\dimexpr0.98\linewidth-#1\linewidth\relax}\raggedright #3\end{minipage}\par}
% Linke Spalte mit Teilaufgabe (Buchstabe hängt links heraus wie in der Liste)
\newcommand{\kbspalte}[1]{\list{}{\leftmargin\z@\labelwidth1.4em\labelsep0.2em\topsep\z@
  \partopsep\z@\parsep\z@\itemsep\z@}\raggedright #1\endlist}
% Antwortfelder: 2,2 cm, in Punkten 1,2 cm; ohne Umbruch davor
\newcommand{\kbfeld}[1][]{\mbox{\underline{\hspace{2.2cm}}\ifblank{#1}{}{\,#1}}}
\newcommand{\kbfeldk}[1][]{\mbox{\underline{\hspace{1.2cm}}\ifblank{#1}{}{\,#1}}}
% Grafik oder Tabelle unter dem Text, linksbündig (Befund 8, 28)
\newcommand{\kbdarunter}[1]{\par\vspace{4pt}\noindent#1\par}
% Ankreuzen: je Option eine Zeile, linksbündig (Befund 6, 7)
\newcommand{\kbkreuz}[1]{\par\vspace{2pt}\noindent$\square$\ #1\par}
% Bearbeitungsraum (Befund 9): Schreibzeilen wie mathblatt (9 mm)
\newcommand{\kbraum}[1]{\schreibzeilen{#1}}
% Wertetabelle mit Köpfen im Textmodus (Befund 26): wie \wertetabelle der
% Vorlage, aber die Köpfe setzt der Aufruf selbst ($x$, „Zeit in h“).
\NewDocumentCommand{\kbwertetabelle}{o m m m}{%
  \def\mbkopf{}\def\mbwerte{}\def\mbleer{}\setcounter{mbn}{0}%
  \foreach \v in {#4}{\xappto\mbkopf{& $\v$}\xappto\mbleer{& }\stepcounter{mbn}\xdef\mbnn{\arabic{mbn}}}%
  \IfNoValueTF{#1}{}{\foreach \v in {#1}{\ifdefempty{\v}{\xappto\mbwerte{& }}{\xappto\mbwerte{& $\v$}}}}%
  \begingroup\renewcommand{\arraystretch}{1.3}%
  \edef\mbpre{|>{\noexpand\columncolor{mbkasten}}l|*{\mbnn}{>{\noexpand\centering\noexpand\arraybackslash}p{\noexpand\mbzell}|}}%
  \expandafter\mbtabstart\expandafter{\mbpre}\hline
  #2 \mbkopf \\ \hline
  \IfNoValueTF{#1}{\rule{0pt}{\mbzeile}#3 \mbleer \\ \hline}{#3 \mbwerte \\ \hline}
  \end{tabular}\endgroup}
% Merkkasten am Ende (Aufbau des Kompetenzblatts)
\newcommand{\kbmerk}[2]{\par\addvspace{14pt}\begin{lrbox}{\kbbox}\begin{minipage}[t]{\linewidth}%
  \noindent{\bfseries #1}\par\vspace{1pt}\noindent{\color{mbgrau}\rule{\linewidth}{0.4pt}}\par\vspace{4pt}%
  \uebersichtskasten{\raggedright #2}\end{minipage}\end{lrbox}%
  \setlength{\kbhoehe}{\dimexpr\ht\kbbox+\dp\kbbox\relax}\Needspace*{\kbhoehe}\noindent\usebox{\kbbox}\par}
% Lösungsblatt: Nummer und Buchstabe halbfett, Lösung daneben
\newcommand{\kbloesung}[2]{\par\addvspace{3pt}\noindent\hangindent2.8em\hangafter1\makebox[2.8em][l]{\textbf{#1}}#2\par}
% Zeichen dieses Blatts, die Latin Modern nicht hat
\newunicodechar{α}{\kbzeichen{α}{\alpha}}
\newunicodechar{β}{\kbzeichen{β}{\beta}}
\newunicodechar{γ}{\kbzeichen{γ}{\gamma}}
\newunicodechar{≈}{\kbzeichen{≈}{\approx}}
\makeatother

\begin{document}
\kbkopfzeile{Trigonometrie · Kompetenzblatt}{TRI-K1}
% Kompetenzblatt TRI-K1: trigonometrie e4 „Sinussatz“ (zusammenbau v0.7, Niveau FOR)
\kbtitel{Ich kann Seiten und Winkel mit dem Sinussatz berechnen.}{ab Kl. 10 (am Gymnasium ab Kl. 9) · P10 ×9}
\setcounter{aufgabe}{0}

\kbabschnitt{Das kennst du schon}

% zone: trigonometrie-zone-f7-v1
\begin{kbaufgabe}{Ich kann den dritten Winkel mit der Winkelsumme berechnen.}
\begin{kbblock}
\kbzweispaltig[0.46]{\kbspalte{\kbteil{}{Ein Dreieck hat die Winkel $45^\circ$ und $85^\circ$.}{}\kbfrage{Wie groß ist der dritte?}\kbantwort{\kbfeld $^\circ$}}}{\kbraum{2}}
\end{kbblock}
\end{kbaufgabe}

% zone: trigonometrie-zone-f4-v1
\begin{kbaufgabe}{Ich kann eine Gleichung mit Bruch umstellen.}
\begin{kbblock}
\kbzweispaltig[0.46]{\kbspalte{\kbanweisung{Löse die Gleichung.}\kbteil{}{{\bfseries\boldmath $\frac{x}{5} = 3$}}{}\kbfrage{Wie groß ist $x$?}\kbantwort{$x$ = \kbfeld}}}{\kbraum{2}}
\end{kbblock}
\end{kbaufgabe}

% zone: trigonometrie-zone-f8-v1
\begin{kbaufgabe}{Ich kann Längen in eine andere Einheit umrechnen.}
\begin{kbblock}
\kbanweisung{Rechne um.}
\kbteil{}{{\bfseries\boldmath $3$ km}}{}
\kbantwort{\kbfeld[m]}
\end{kbblock}
\end{kbaufgabe}

\kbabschnitt{Schritt für Schritt}

% leiter: trigonometrie-e4-k1-s0-v1
\begin{kbaufgabe}{Ich kann erkennen, ob ich mit sin, cos, tan oder mit dem Sinussatz rechne.}
\begin{kbblock}
\kbteil{}{Im Dreieck $ABC$ ist bei $B$ ein rechter Winkel markiert, $a$ und $\alpha$ sind gegeben.}{}
\kbfrage{Womit rechnest du $b$?}
\kbkreuz{sin, cos, tan}
\kbkreuz{Sinussatz}
\end{kbblock}
\end{kbaufgabe}

% leiter: trigonometrie-e4-k1-s1-v1, trigonometrie-e4-k1-s2-v1 (Paar, dicht)
\kbpaar{
\kbnummer{Ich kann eine Seite mit dem Sinussatz berechnen.}\kbhalb{\kbteil{}{Dreieck $ABC$: $a = 7$ cm, $\alpha = 68^\circ$, $\beta = 51^\circ$.}{}
\kbfrage{Wie lang ist $b$?}
\kbantwort{$b$ ≈ \kbfeld[cm]}
\item[]\kbraum{2}}
}{
\kbnummer{Ich kann Seiten und Gegenwinkel auch bei anderen Eckennamen zuordnen.}\kbhalb{\kbteil{}{Dreieck $PQR$: $PQ = 8$ cm, Winkel bei $R$: $69^\circ$, Winkel bei $Q$: $53^\circ$.}{}
\kbfrage{Wie lang ist $PR$?}
\kbantwort{$PR$ ≈ \kbfeld[cm]}
\item[]\kbraum{2}}
}

% leiter: trigonometrie-e4-k1-s3-v1, trigonometrie-e4-k1-s4-v2 (Paar, dicht)
\kbpaar{
\kbnummer{Ich kann den fehlenden Winkel erst mit der Winkelsumme berechnen.}\kbhalb{\kbteil{}{Dreieck $ABC$: $a = 8$ cm, $\alpha = 63^\circ$, $\beta = 48^\circ$.}{}
\kbfrage{Wie lang ist $c$?}
\kbantwort{$c$ ≈ \kbfeld[cm]}
\item[]\kbraum{2}}
}{
\kbnummer{Ich kann den Sinussatz mit einem stumpfen Winkel anwenden.}\kbhalb{\kbteil{}{Dreieck $ABC$: $b = 9$ cm, $\beta = 127^\circ$, $\gamma = 24^\circ$.}{}
\kbfrage{Wie lang ist $c$?}
\kbantwort{$c$ ≈ \kbfeld[cm]}
\item[]\kbraum{2}}
}

% leiter: trigonometrie-e4-k1-s5-v1, trigonometrie-e4-k1-s6-v1 (Paar, dicht)
\kbpaar{
\kbnummer{Ich kann einen stumpfen Winkel in der Winkelsumme verwenden.}\kbhalb{\kbteil{}{Dreieck $ABC$: $a = 7$ cm, $\alpha = 27^\circ$, $\beta = 118^\circ$.}{}
\kbfrage{Wie lang ist $c$?}
\kbantwort{$c$ ≈ \kbfeld[cm]}
\item[]\kbraum{2}}
}{
\kbnummer{Ich kann mit Dezimalzahlen und Einheiten rechnen und runden.}\kbhalb{\kbteil{}{Dreieck $ABC$: $a = 4{,}35$ m, $\alpha = 58{,}2^\circ$, $\beta = 47{,}6^\circ$.}{}
\kbfrage{Wie lang ist $b$, auf eine Dezimale?}
\kbantwort{$b$ ≈ \kbfeld[m]}
\item[]\kbraum{2}}
}

% leiter: trigonometrie-e4-k1-s7-v1, trigonometrie-e4-k1-s8-v1 (Paar, dicht)
\kbpaar{
\kbnummer{Ich kann ein vorgegebenes Ergebnis mit einer Rechnung zeigen.}\kbhalb{\kbteil{}{Dreieck $ABC$: $a = 9$ cm, $\alpha = 77^\circ$, $\beta = 44^\circ$.}{}
\kbfrage{Zeige rechnerisch, dass $b \approx 6{,}4$ cm ist.}
\item[]\kbraum{2}}
}{
\kbnummer{Ich kann zwei Seiten berechnen und eine Aussage prüfen.}\kbhalb{\kbteil{}{Dreieck $ABC$: $c = 10$ cm, $\alpha = 53^\circ$, $\beta = 56^\circ$.}{}
\kbfrage{Prüfe die Aussage „$a$ ist länger als $b$.“}
\item[]\kbraum{2}}
}

% leiter: trigonometrie-e4-k1-s9-v1, trigonometrie-e4-k1-s10-v2 (Paar, dicht)
\kbpaar{
\kbnummer{Ich kann den Sinussatz im Viereck anwenden.}\kbhalb{\kbteil{}{Viereck $ABCD$ mit rechtem Winkel bei $B$ und der Diagonale $AC = 8$ m. $\angle BAC = 57^\circ$, $\angle BCD = 91^\circ$, $\angle ADC = 76^\circ$.}{}
\kbfrage{Wie lang ist $AD$?}
\kbantwort{$AD$ ≈ \kbfeld[m]}
\item[]\kbraum{2}}
}{
\kbnummer{Ich kann mit dem Ergebnis eine Weglänge weiterrechnen.}\kbhalb{\kbteil{}{Dreieck $ABC$: $b = 16$ m, $\beta = 104^\circ$, $\alpha = 31^\circ$. Der Punkt $P$ liegt auf $BC$, $PC = 5$ m.}{}
\kbfrage{Wie lang ist $BP$?}
\kbantwort{$BP$ ≈ \kbfeld[m]}
\item[]\kbraum{2}}
}

% leiter: trigonometrie-e4-k1-s11-v1, trigonometrie-e4-k1-s12-v1 (Paar, dicht)
\kbpaar{
\kbnummer{Ich kann ein Ergebnis über die Höhe nachrechnen.}\kbhalb{\kbteil{}{Dreieck $ABC$: $a = 8$ cm, $\alpha = 67^\circ$, $\beta = 53^\circ$.}{}
\kbfrage{Berechne $b$ mit dem Sinussatz und noch einmal über die Höhe $h$ von $C$ auf $AB$.}
\kbantwort{$b$ ≈ \kbfeld[cm]}
\item[]\kbraum{2}}
}{
\kbnummer{Ich kann einen Winkel mit dem Sinussatz berechnen.}\kbhalb{\kbteil{}{Dreieck $ABC$: $a = 9$ cm, $\alpha = 68^\circ$, $b = 7$ cm.}{}
\kbfrage{Wie groß ist $\beta$?}
\kbantwort{$\beta \approx$ \kbfeld $^\circ$}
\item[]\kbraum{2}}
}

% leiter: trigonometrie-e4-k1-s13-v1, trigonometrie-e4-k1-s14-v1 (Paar, dicht)
\kbpaar{
\kbnummer{Ich kann die dritte Seite mit dem Kosinussatz berechnen.}\kbhalb{\kbteil{}{Dreieck $ABC$: $a = 5$ cm, $b = 8$ cm, $\gamma = 47^\circ$.}{}
\kbfrage{Wie lang ist $c$?}
\kbantwort{\kbfeld[cm]}
\item[]\kbraum{2}}
}{
\kbnummer{Ich kann einen Winkel mit dem Kosinussatz berechnen.}\kbhalb{\kbteil{}{Dreieck $ABC$: $a = 5$ cm, $b = 7$ cm, $c = 6$ cm.}{}
\kbfrage{Wie groß ist $\gamma$?}
\kbantwort{$\gamma \approx$ \kbfeld $^\circ$}
\item[]\kbraum{2}}
}

% leiter: trigonometrie-e4-k1-s15-v1
\begin{kbaufgabe}{Ich kann den Sinussatz über die Höhe herleiten.}
\begin{kbblock}
\kbteil{}{Die Höhe $h_c$ zerlegt das Dreieck $ABC$ in zwei rechtwinklige Teildreiecke.}{}
\kbfrage{Drücke $h_c$ in beiden Teildreiecken aus und leite daraus $\frac{a}{\mathrm{sin}\,\alpha} = \frac{b}{\mathrm{sin}\,\beta}$ her.}
\item[]\kbraum{2}
\end{kbblock}
\end{kbaufgabe}

% pruefung: trigonometrie-e4-k1-s16-v3, trigonometrie-e4-k1-s16-v1, trigonometrie-e4-k1-s16-v11, trigonometrie-e4-k1-s16-v9, trigonometrie-e4-k1-s16-v7
\begin{kbaufgabe}{Ich kann das auch in Aufgaben aus der Prüfung.}
\begin{kbblock}
\kbteil{a)}{Eine Rampe $y$ führt auf eine Treppe. Im Dreieck aus Rampe, waagerechter Strecke ($220$ cm) und der Verbindung Treppenfuß–obere Stufenkante liegt am Rampenfuß der Winkel $6^\circ$ und am Treppenfuß der Winkel $137^\circ$ (gegenüber $y$).}{P10 ’25}
\kbfrage{Berechne $y$.}
\kbzweispaltig[0.46]{\kbantwort{$y$ ≈ \kbfeld[cm]}}{\kbraum{3}}
\end{kbblock}
\begin{kbblock}
\kbteil{b)}{Eine Seilbahn führt von $A$ über die Stütze $B$ nach $C$. $AB = 520$ m, der Winkel bei $A$ ist $31^\circ$, der Winkel $ABC$ ist $113^\circ$.}{P10 ’24}
\kbfrage{Berechne die Länge der Strecke $BC$.}
\kbzweispaltig[0.46]{\kbantwort{$BC$ ≈ \kbfeld[m]}}{\kbraum{3}}
\end{kbblock}
\begin{kbblock}
\kbteil{c)}{Viereck $ABCD$ mit rechtem Winkel bei $A$ und der Diagonale $BD$. $\angle ADB = 48^\circ$, $\angle BDC = 57^\circ$, $\angle ABC = 104^\circ$, $DC = 1\,300$ m.}{P10 ’21}
\kbfrage{Stimmt die Aussage „$BC$ ist genauso lang wie $DC$“?}
\item[]\kbraum{3}
\end{kbblock}
\begin{kbblock}
\kbteil{d)}{Dreieck $DEF$: $FE = 9$ m, Winkel bei $F$ $118^\circ$, Winkel bei $E$ $23^\circ$. $P$ liegt auf $DE$ mit $PE = 8$ m.}{P10 ’20}
\kbfrage{Wie lang ist $DP$?}
\kbzweispaltig[0.46]{\kbantwort{$DP$ ≈ \kbfeld[m]}}{\kbraum{3}}
\end{kbblock}
\begin{kbblock}
\kbteil{e)}{Ein Grundstück $ABCD$ hat bei $B$ einen rechten Winkel. $AC = 7$ m, $\angle BAC = 59^\circ$, $\angle BCD = 88^\circ$, $\angle ADC = 78^\circ$.}{P10 ’19}
\kbfrage{Berechne den Abstand $AD$.}
\kbzweispaltig[0.46]{\kbantwort{$AD$ ≈ \kbfeld[m]}}{\kbraum{3}}
\end{kbblock}
\end{kbaufgabe}

\kbmerk{Zum Merken}{Sinussatz: In jedem Dreieck ist das Verhältnis aus einer Seite und dem Sinus ihres Gegenwinkels für alle drei Seiten gleich – a/sin α = b/sin β = c/sin γ. Man braucht ein vollständiges Paar (Seite und Gegenwinkel) und von einem zweiten Paar eine Größe. \\ b = 8 cm, β = 40°, α = 75°: a = 8 · sin 75° : sin 40° ≈ 12,0 cm \\ Fehlt der Winkel gegenüber der gesuchten Seite, liefert ihn die Winkelsumme: γ = 180° − 75° − 40° = 65°. \\ Ein stumpfer Winkel bleibt, wie er ist – der Taschenrechner rechnet sin 120° direkt; in die Winkelsumme gehört der stumpfe Winkel selbst, nicht sein Nebenwinkel.}
\end{document}
